\documentclass[UTF8]{ctexart}
\usepackage{amsmath,amssymb}

\title{BMS(Bashicu Matrix System) 多版本计算规则}
\author{no-name5844}

\begin{document}
\maketitle

\setcounter{section}{-1}

BMS 是 Bashicu 于 2014 年提出的序数 / 巨大数记号，记法为按列书写的非负整数矩阵，通过对末列做\textbf{展开}递归生成函数。所有版本共用同一套展开骨架，唯一差异在\textbf{父项条件}与\textbf{上升度数列 $a_{k,m}$}（决定副本中是否加 $a\cdot\Delta_k$）。本笔记先给以 BM4 为准的共用骨架，再逐版本列出其上升度 / 父项差异。格式仿照序列记号笔记（定义 $\to$ 展开 $\to$ 辅助函数 $\to$ expand 函数）。

记号：列标 $x$、行标 $y$（从 $0$ 起）。$S=S_0S_1\cdots S_{X-1}$（$X$ 列）；列 $S_x=(S_{x,0},\dots,S_{x,Y-1})$。$S_{x,y}$ 为第 $x$ 列第 $y$ 行元素。$S+T$ 为列拼接。末列 $S_{X-1}$；其最靠下非零元素所在行 $z=\max\{y\mid S_{X-1,y}>0\}$（该元素称 LNZ）。
满足各列之间的字典序比较，在两个列之间的大小为比较为各元素的字典（下标为从小到大）。

\section{共用展开骨架（BM4 标准；其余版本共用）}
\subsection{定义}
合法矩阵（不一定标准型）满足：
\begin{enumerate}
  \item 首列全 $0$：$S_{0,y}=0$；
  \item 每列自上而下非增：$S_{x,y}\ge S_{x,y+1}$；
  \item 同行不超前超 $1$：$S_{x,y}\le\max_{p<x}S_{p,y}+1$（空最大值为 $0$）。
\end{enumerate}
极限表达式为 $,(0),(0)(1),(0)(1,1),(0)(1,1,1),\dots$

\subsection{展开}
\subsubsection{辅助函数}
\begin{enumerate}
  \item \textbf{父项} $p_k(m)$：取满足
  \[
  S_{p,k}<S_{m,k}\quad\text{且（}k>0\text{ 时）}\quad a_{k-1,m}(p)=1
  \]
  的最大 $p<m$。其中 $a_{k-1,m}(p)=1\iff\exists c\ge0:(p_{k-1})^c(m)=p$（列 $p$ 是列 $m$ 在 $k-1$ 行下的祖先）。约定 $(p_k)^0(m)=m$。
  \item \textbf{上升度} $a_{k,m}(t)$：对坏部第 $m$ 列（$t$ 为坏根列标），
  \[
  a_{k,m}(t)=\begin{cases}1,& \exists c\ge0:(p_k)^c(m)=t\\0,& \text{否则}\end{cases}
  \]
  \item \textbf{坏根} $r=p_z(X-1)$（末列 LNZ 的父项）；末列全 $0$ 则无坏根。
  \item \textbf{好部 / 坏部}：$G=S_0+\cdots+S_{r-1}$，$B=S_r+\cdots+S_{X-2}$。
  \item \textbf{阶差} $\Delta_k=\begin{cases}S_{X-1,k}-S_{r,k},& k<z\\0,& k\ge z\end{cases}$。
  \item \textbf{副本} $B^{(i)}$ 第 $m$ 列：$B^{(i)}_{m,k}=B_{m,k}+a_{k,r+m}(r)\cdot\Delta_k\cdot i$。
\end{enumerate}

\subsubsection{expand 函数}
\begin{align*}
\mathit{expand}((),n)&=n\\
\mathit{expand}(S+(0),n)&=S\\
\mathit{expand}(S,n)&=G+B^{(0)}+\cdots+B^{(n)}\quad(\text{末列非全 }0)
\end{align*}
基本列前驱 $FS_n(S)=G+B^{(0)}+\cdots+B^{(n-1)}$。序数映射：
\[
\begin{aligned}
&\mathit{Trans}()=0,\quad\mathit{Trans}(S+(0,\dots,0))=\mathit{Trans}(S)+1,\\
&\mathit{Trans}(M)=\sup\{\mathit{Trans}(FS_i(M))\mid i\in\mathbb N\}.
\end{aligned}
\]

\section{BM1(初版, 2014)}
\subsection{定义}
同 \S 0.1。

\subsection{展开}
\subsubsection{辅助函数}
父项 $p_k(m)$（$k>0$）\textbf{缺上行祖先检查}（即 $a_{k-1,m}(p)=1$ 条件不生效），退化为仅按首行判定。上升度由首行统一决定：
\[
a_{k,m}(t)=a_{0,m}(t)\quad(\forall k)
\]
（各行的祖先链判定被忽略，只用首行父项链。）

\subsubsection{expand 函数}
同 \S 0.2.2。
终止性：否。反例：讨论页给出的 non-terminating sequence（无具体矩阵式公开）。

\section{BM2(第二版, 2016)}
\subsection{定义}
同 \S 0.1。

\subsection{展开}
\subsubsection{辅助函数}
引入逐行上升度，但父项 / 上升度判定规则与 BM2.3 不同（坏部某些列的祖先链处理有误）；其精确逐条公式未在公开文档中给出，差异落在 $a_{k,m}$ 的计算。

\subsubsection{expand 函数}
同 \S 0.2.2。
终止性：否。反例矩阵（与 BM2.3 展开不同且无限循环）：
\[
\begin{aligned}
&(0,0,0,0)(1,1,1,1)(2,2,1,1)(3,3,1,1)\\
&(4,2,0,0)(5,1,1,1)(6,2,1,1)(7,3,1,1).
\end{aligned}
\]

\section{BM2.1(koteitan 过渡版)}
\subsection{定义}
同 \S 0.1。

\subsection{展开}
\subsubsection{辅助函数}
BM2 $\to$ BM2.3 之间的中间修正；具体 $a_{k,m}$ 规则见 basmat 计算器源码（按 \texttt{version} 分支），公开文档未给逐条公式。

\subsubsection{expand 函数}
同 \S 0.2.2。终止性：未单独证明（接近 BM2.3，故同 BM4）。

\section{BM2.2(koteitan 过渡版)}
\subsection{定义}
同 \S 0.1。

\subsection{展开}
\subsubsection{辅助函数}
同 BM2.1：BM2 $\to$ BM2.3 的过渡修正，精确 $a_{k,m}$ 规则见 basmat 源码。

\subsubsection{expand 函数}
同 \S 0.2.2。终止性：未单独证明。

\section{BM2.3(koteitan 修补版, 2018)}
\subsection{定义}
同 \S 0.1。

\subsection{展开}
\subsubsection{辅助函数}
采用 \S 0.2.1 的父项与上升度规则（含上行祖先检查、逐行独立判定）。\textbf{即与 BM4 完全相同}。

\subsubsection{expand 函数}
同 \S 0.2.2（= BM4）。终止性：同 BM4。

\section{BM3(第三版, 2018-06-12)}
\subsection{定义}
同 \S 0.1。

\subsection{展开}
\subsubsection{辅助函数}
逐行上升度（比 BM1 接近 BM4），但 $k>0$ 行的父项 $p_k$ 条件与 BM4 \textbf{不同}（不完整：祖先链检查与 BM4 不一致），故 $a_{k,m}$ 在坏部某些列上误判。

\subsubsection{expand 函数}
同 \S 0.2.2。
终止性：否。反例矩阵（社区记为无限循环示意）：
\[
(0,0,0)(1,1,1)(2,1,1)(1,1,1).
\]

\section{BM3.1(Nish 变体)}
\subsection{定义}
同 \S 0.1。

\subsection{展开}
\subsubsection{辅助函数}
BM3 的社区修正；精确 $a_{k,m}$ 规则见 basmat 源码。已知行为：在 $(0,0,0)\allowbreak(1,1,1)\allowbreak(2,1,0)\allowbreak(1,1,1)$ 上展开整齐，但在 $(0,0,0)\allowbreak(1,1,1)\allowbreak(2,2,0)\allowbreak(1,1,1)$ 上该性质失效。

\subsubsection{expand 函数}
同 \S 0.2.2。终止性：未证明（有奇异行为）。

\section{BM3.2(Nish 变体)}
\subsection{定义}
同 \S 0.1。

\subsection{展开}
\subsubsection{辅助函数}
同 BM3.1：BM3 的进一步修正，精确 $a_{k,m}$ 规则见 basmat 源码。行为同上（在 $(0,0,0)(1,1,1)(2,2,0)(1,1,1)$ 失效）。

\subsubsection{expand 函数}
同 \S 0.2.2。终止性：未证明（奇异行为未消除）。

\section{BM3.3(Rpakr + ecl1psed, 2019-03)}
\subsection{定义}
同 \S 0.1。

\subsection{展开}
\subsubsection{辅助函数}
与 BM4 唯一差异在上升度。设 $t=$ LNZ 行、坏根列 $r$。坏部各列分两类：
\begin{itemize}
  \item \textbf{第 $t$ 行是坏根 $r$ 的后代（含 $r$ 本身）的列}：前 $t-1$ 行 $a_{k,m}$ 全为 $1$（与 BM4 相同）。
  \item \textbf{第 $t$ 行不是坏根后代的列}：
  \[
  \text{BM4}:\quad a_{k,m}=1\iff a_{k,\mathrm{parent}(m)}=1\ \land\ (\mathrm{parent}=r\ \lor\ \mathrm{parent}>r)
  \]
  \[
  \text{BM3.3}:\quad a_{k,m}=1\iff a_{k,\mathrm{parent}(m)}=1\ \land\ \mathrm{parent}>r
  \]
  （BM3.3 去掉了“父项 = 坏根”这一条件。）
\end{itemize}

\subsubsection{expand 函数}
同 \S 0.2.2（仅上升度按 9.2.1）。终止性：未证明。
示例：$(0,0,0)\allowbreak(1,1,1)\allowbreak(2,1,0)\allowbreak(1,1,1)$ 在 BM3.3 下展开为
\[
\begin{aligned}
&(0,0,0)(1,1,1)(2,1,0)(1,1,0)(2,2,1)(3,1,0)\\
&(2,2,0)(3,3,1)(4,1,0)(3,3,0)(4,4,1)(5,1,0)\dots
\end{aligned}
\]

\section{BM4(现行标准, 2018-09-01)}
BM4 = BM2.3 标准化，展开规则即 \S 0 骨架。
\subsection{定义}
同 \S 0.1。

\subsection{展开}
\subsubsection{辅助函数}
同 \S 0.2.1（父项含上行祖先检查；上升度按各行独立判定祖先链）。

\subsubsection{expand 函数}
同 \S 0.2.2。
Bashicu 命名的对应序数：$(0,0,0)\allowbreak(1,1,1)\allowbreak(2,2,0)\to$ first back gear；$(0,0,0)\allowbreak(1,1,1)\allowbreak(2,2,1)\allowbreak(3,3,0)\to$ second back gear；$(0,0,0)\allowbreak(1,1,1)\allowbreak(2,2,2)\to\omega$ back ordinal。
终止性：声称良基（Yto / Rachel Hunter, arXiv:2307.04606, 2023--2024），未评审。

\section{PsiCubed2 版本(社区个人变体)}
\subsection{定义}
同 \S 0.1。

\subsection{展开}
\subsubsection{辅助函数}
个人变体，具体 $a_{k,m}$ 规则见其博客（User blog:PsiCubed2/My own version of BMS）；差异同样落在父项 / 上升度。

\subsubsection{expand 函数}
同 \S 0.2.2。终止性：未评估。

\section{Idealized BMS 系列(社区提案)}
目标：使 $A(1,1,1)\leftrightarrow\psi(C+\Omega_\omega)$ 成立（更整齐的序数对应）。各提案仅在上升度 / 坏部判定 / 父项检查上微调，均未被采纳为标准。提案名：
\texttt{Mar.19(1)}, \texttt{Mar.19(2)}, \texttt{BR+delta comp}, \texttt{UBRABC}, \texttt{UBRABC+parent check}, \texttt{UBRA+parent check}, \texttt{Rpakr Def.1}--\texttt{Def.5}.
计算规则见各提案原文；终止性均为开放问题。

\end{document}
